%% kapa: exam=1 round=22 year=2011 period=2 form=A subject=accounting subject_ko=회계학 seq=34 number=34
%% source: Q-Net 감정평가사 기출문제 (cst00309) artlSeq=5036790 file=extracted/1008/5036790/2011년 제22회 감정평가사 1차시험 기출문제/2교시.hwp
%% conversion: kordoc 4.24.51 -> tools/gen1.py (auto); review=proofread-sonnet; flags=none
\begin{question}{34}
(주)서울은 제1공정에서 제품A와 B를 생산하고 있으며 20×1년 생산량 및 원가자료는 다음과 같다.

{\small\centering
\begin{tabular}{|l|r|r|r|}
\hline
 & \multicolumn{1}{c|}{제품A} & \multicolumn{1}{c|}{제품B} & \multicolumn{1}{c|}{계} \\
\hline
결합원가 & ? & ? & ￦50,000 \\
생산량 & 200단위 & 300단위 & \\
판매단가 & ￦2,000 & ￦2,500 & \\
\hline
\end{tabular}\par}

제품A는 제2공정에서 제품C 200단위(판매단가 ￦3,000)로 추가 가공할 수 있다. 제품A를 추가 가공하는데 소요된 원가는 ￦200,000이다. 제품A의 판매비는 ￦150,000이고, 제품C의 판매비는 ￦100,000이다. 제품A를 모두 제품C로 추가 가공하여 판매하는 경우 이익은 얼마나 증가하는가? \par
\begin{choices}
  \item ￦10,000
  \item ￦20,000
  \item ￦30,000
  \item ￦40,000
  \item ￦50,000
\end{choices}
\end{question}
\answer{5}
